\chapter{Opérations \& débogage}
\label{ch:operations-debugging}

Ce chapitre documente l'emplacement des journaux sur disque, leur
niveau de verbosité configurable, et les commandes courantes utiles
quand un utilisateur signale un comportement étrange et que vous
voulez voir ce que l'application a réellement fait.

\section{Où sont les journaux}
\label{sec:logs-location}

Depuis la \textbf{v1.7.0}, mybibli écrit ses journaux à deux endroits
simultanément :

\begin{itemize}
  \item \textbf{stdout} (inchangé) — accessible via
        \texttt{docker compose logs mybibli}. Idéal pour suivre le
        flux en direct pendant une session active.
  \item \textbf{Un fichier rotatif quotidien sur disque}, par défaut
        sous \texttt{/var/log/mybibli/} dans le conteneur, monté sur
        un volume nommé (\texttt{mybibli\_logs}) côté hôte pour que
        les fichiers survivent à \texttt{docker compose pull \&\&
        up -d}. Idéal pour le débogage \emph{post-mortem}.
\end{itemize}

Les fichiers sont nommés \texttt{mybibli.log.AAAA-MM-JJ} (un par jour
UTC). Une tâche en arrière-plan purge automatiquement les fichiers de
plus de 30 jours.

\subsection*{Changer le chemin des journaux dans le conteneur}
\label{sec:log-dir-override}

Le répertoire de journaux par défaut \texttt{/var/log/mybibli/} dans
le conteneur est contrôlé par la variable d'environnement
\texttt{MYBIBLI\_LOG\_DIR}. Ne le changez que si vous remontez le
volume ailleurs — le fichier \texttt{docker-compose.yml} fourni
mappe le volume nommé \texttt{mybibli\_logs} (ou un bind mount via
\texttt{LOGS\_HOST\_PATH}) sur \texttt{/var/log/mybibli}, et changer
l'un sans l'autre fait atterrir les journaux dans la couche
inscriptible où ils sont effacés à chaque
\texttt{docker compose up -d} suivant un \texttt{pull}.

Voir section~\ref{sec:install-logs-bind} pour la configuration du
bind mount (journaux visibles côté hôte pour DSM Log Center /
journald shipping).

\section{Lire les journaux depuis l'hôte}

\begin{lstlisting}
# Suivre en temps réel le fichier du jour
docker compose exec mybibli tail -f /var/log/mybibli/mybibli.log.$(date -u +%Y-%m-%d)

# Dernières 100 lignes du fichier du jour
docker compose exec mybibli tail -n 100 /var/log/mybibli/mybibli.log.$(date -u +%Y-%m-%d)

# Lister tous les jours disponibles
docker compose exec mybibli ls /var/log/mybibli/

# Rechercher un mot-clé sur tous les jours conservés
docker compose exec mybibli sh -c "grep ERROR /var/log/mybibli/mybibli.log.*"
\end{lstlisting}

Si vous avez monté le volume comme un bind mount
(\texttt{LOGS\_HOST\_PATH} dans \texttt{.env}, voir
section~\ref{sec:install-logs-bind}), les fichiers sont visibles
directement depuis l'hôte sans passer par \texttt{docker compose
exec}.

\section{Niveaux de journalisation}
\label{sec:log-levels}

Deux variables d'environnement contrôlent la verbosité :

\begin{itemize}
  \item \texttt{MYBIBLI\_LOG\_LEVEL} (recommandée, v1.7.0+) —
        accepte soit un niveau simple (\texttt{trace}, \texttt{debug},
        \texttt{info}, \texttt{warn}, \texttt{error}) soit une liste
        de directives \texttt{tracing-subscriber}
        \texttt{EnvFilter}, par exemple
        \texttt{mybibli=debug,sqlx::query=warn}.
  \item \texttt{RUST\_LOG} — fallback historique, même format.
        Pris en compte si \texttt{MYBIBLI\_LOG\_LEVEL} est absent.
\end{itemize}

La valeur par défaut est \texttt{info} — sûre en production (info,
warn, error ; pas de bruit SQL par requête, pas de flot de traces
par requête HTTP).

\subsection*{Niveaux recommandés}

\begin{description}
  \item[\texttt{info}] Défaut production. Capture les événements
        d'authentification, les récupérations de métadonnées, les
        tâches planifiées, les erreurs visibles utilisateur. Environ
        10 à 100 lignes par jour pour une instance familiale.
  \item[\texttt{mybibli=debug}] Investigation active d'une
        fonctionnalité mybibli. Ajoute le routage des requêtes, les
        cache hits/misses, les décisions du dispatcher de scan.
        Environ 10$\times$ le volume \texttt{info}.
  \item[\texttt{mybibli=trace}] Investigation profonde ponctuelle.
        Inclut le tracing par appel dans le code mybibli. Très
        verbeux ; à utiliser brièvement puis revenir à
        \texttt{info}.
  \item[\texttt{mybibli=debug,sqlx::query=info}] Investigation active
        ET voir chaque instruction SQL (utile quand une requête est
        suspecte).
\end{description}

\subsection*{Comment changer le niveau}

\textbf{Deux chemins}, selon que vous voulez que le changement
survive aux redémarrages du conteneur ou qu'il ne dure que le
temps d'une investigation.

\subsubsection*{Depuis l'interface admin (v1.7.1+, recommandé)}

Connectez-vous en Admin, allez sur \texttt{/admin?tab=system},
descendez jusqu'à la section \textbf{Journalisation}. Saisissez la
nouvelle directive dans le champ \texttt{Niveau de journalisation}
(tout ce que \texttt{tracing-subscriber::EnvFilter} accepte
\string— niveau simple ou liste de directives). Cliquez
\textbf{Enregistrer}. La prochaine ligne de log utilise le nouveau
filtre \string— sans redémarrage du conteneur, sans
\texttt{docker compose up -d}. La valeur est persistée en table
\texttt{settings} et survit aux redémarrages.

\subsubsection*{Depuis \texttt{docker-compose.yml} / \texttt{.env}}

Modifier le fichier, fixer \texttt{MYBIBLI\_LOG\_LEVEL}, puis :

\begin{lstlisting}
docker compose up -d
\end{lstlisting}

Le conteneur redémarre et le nouveau niveau prend effet dès la
prochaine ligne de journal. \textbf{À noter :} la valeur enregistrée
en admin dans la table \texttt{settings} gagne au prochain boot.
Donc si vous avez déjà enregistré un niveau depuis l'UI admin, c'est
celui-là qui sera chargé \string— la variable d'environnement ne
compte qu'à l'installation fraîche avant le premier enregistrement.

\section{Lire le JSON structuré}

Les journaux sont émis en JSON ligne par ligne. Chaque ligne porte :

\begin{itemize}
  \item \texttt{timestamp} — ISO~8601 UTC.
  \item \texttt{level} — \texttt{TRACE} / \texttt{DEBUG} /
        \texttt{INFO} / \texttt{WARN} / \texttt{ERROR}.
  \item \texttt{target} — le module Rust qui a émis la ligne
        (par ex. \texttt{mybibli::routes::catalog}). Filtrer sur ce
        champ pour restreindre une recherche à un sous-système.
  \item \texttt{fields.message} — le résumé lisible.
  \item d'autres \texttt{fields.*} — clés de contexte comme
        \texttt{user\_id}, \texttt{volume\_id}, \texttt{isbn},
        \texttt{title\_id}.
\end{itemize}

Pour un affichage lisible :

\begin{lstlisting}
docker compose exec mybibli tail -n 50 /var/log/mybibli/mybibli.log.$(date -u +%Y-%m-%d) \
  | jq -r '[.timestamp, .level, .fields.message] | @tsv'
\end{lstlisting}

\section{Joindre un extrait de log à un rapport de bug}

Lors de l'ouverture d'une issue sur
\url{https://github.com/guycorbaz/mybibli/issues} :

\begin{enumerate}
  \item Notez l'heure murale où le bug s'est produit (par ex.
        ``autour de 14h30 UTC aujourd'hui'').
  \item Extrayez la fenêtre de 5 minutes concernée :
  \begin{lstlisting}
  # Remplacer 14:25 / 14:35 par votre fenêtre.
  docker compose exec mybibli sh -c \
    "awk '/\"timestamp\":\"$(date -u +%Y-%m-%d)T14:2[5-9]/,/\"timestamp\":\"$(date -u +%Y-%m-%d)T14:3[0-5]/' \
       /var/log/mybibli/mybibli.log.$(date -u +%Y-%m-%d)"
  \end{lstlisting}
  \item Copiez le résultat dans le corps de l'issue, balisé en
        \texttt{```json}.
\end{enumerate}

Caviardez ce que vous estimez sensible — les lignes de journal
peuvent contenir des noms d'emprunteurs, des ISBN de titres que vous
possédez, ou des préfixes de jeton de session (mais jamais les jetons
de session complets ni les hashs de mot de passe ; ils sont
intentionnellement exclus de l'émission de journaux).

\section{Rétention et rotation}

\begin{itemize}
  \item La rotation se fait \textbf{automatiquement} à minuit UTC,
        gérée en interne par \texttt{tracing-appender::rolling}.
  \item Les fichiers de plus de \textbf{30 jours} sont purgés une
        fois par jour par une tâche en arrière-plan. La fenêtre de
        rétention est codée en dur dans la v1.7.0 ; une version
        future l'exposera comme paramètre admin.
  \item Le fichier du jour n'est jamais supprimé (la coupure est
        stricte).
\end{itemize}

Pour conserver les journaux plus de 30 jours, copiez les fichiers
rotatifs hors du volume vers du stockage long terme :

\begin{lstlisting}
docker compose cp mybibli:/var/log/mybibli/. ./logs-archive/
\end{lstlisting}

\section{Réinitialiser un mot de passe administrateur oublié}\label{sec:reset-admin-password}

mybibli ne propose pas de procédure « mot de passe oublié » par
courriel : une installation domestique mono-locataire n'a ni serveur de
messagerie sur lequel s'appuyer, ni second facteur de repli. Si vous
êtes encore connecté quelque part, la correction est ordinaire —
\texttt{/admin > Utilisateurs}, modifiez votre propre compte, donnez-vous
un nouveau mot de passe. Faites-le \emph{avant} l'expiration de cette
session : elle est la seule chose qui vous sépare de la procédure
ci-dessous.

\subsection*{Si vous êtes verrouillé dehors}

Les mots de passe sont stockés sous forme de hachages Argon2id au format
PHC. Les paramètres sont ceux par défaut de la bibliothèque
\texttt{argon2} (\texttt{m=19456,t=2,p=1}), et ils doivent correspondre
exactement, faute de quoi le nouveau mot de passe est rejeté à la
connexion sans aucune explication. \textbf{Ne produisez pas le hachage
avec un outil tiers} : la plupart emploient des paramètres différents.
Utilisez l'utilitaire fourni avec les sources, qui appelle la fonction
même de l'application :

\begin{lstlisting}
git clone https://github.com/guycorbaz/mybibli && cd mybibli
cargo run -q --example hash_password -- 'votre nouveau mot de passe'
\end{lstlisting}

Il affiche une ligne commençant par \texttt{\$argon2id\$v=19\$m=19456}.
Appliquez-la sur la machine qui héberge la base, en remplaçant
\texttt{<HACHAGE>} par cette ligne et \texttt{<UTILISATEUR>} par votre
identifiant :

\begin{lstlisting}
docker compose exec db mariadb -u root -p mybibli \
  -e "UPDATE users SET password_hash = '<HACHAGE>', version = version + 1 \
      WHERE username = '<UTILISATEUR>' AND deleted_at IS NULL;"
\end{lstlisting}

Vérifiez ensuite la ligne avant de refermer la session :

\begin{lstlisting}
docker compose exec db mariadb -u root -p mybibli \
  -e "SELECT username, role, active, deleted_at, \
      LEFT(password_hash,22) FROM users;"
\end{lstlisting}

\subsection*{Trois pièges}

\begin{itemize}
  \item \textbf{Entourez le hachage de guillemets simples.} Il contient
        des caractères \texttt{\$}. Des guillemets doubles laisseraient
        le shell les interpréter et vous écririez une valeur tronquée,
        qui ne pourra jamais correspondre.
  \item \textbf{Conservez \texttt{version = version + 1}.} Le
        verrouillage optimiste lit cette colonne ; l'omettre laisserait
        une modification ultérieure faite depuis l'interface écraser la
        vôtre en croyant partir d'un état à jour.
  \item \textbf{Le mot de passe passe dans l'historique du shell} aux
        deux étapes. Préfixez les commandes d'une espace si votre shell
        est configuré pour ignorer ces lignes, ou purgez l'historique
        ensuite.
\end{itemize}

Il s'agit d'une écriture directe sur des données de production, sans
retour arrière. Lancez d'abord le \texttt{SELECT} : l'opération n'est
réversible que si vous avez noté l'ancienne valeur.

\subsection*{Lire le journal de connexion}

La distinction consignée dans le journal compte pour le diagnostic :
\texttt{Login failed: user not found} signifie qu'aucun compte ne porte
cet identifiant (ou qu'il est supprimé logiquement), tandis que
\texttt{Login failed: invalid password} signifie que le compte est sain
et que seul le mot de passe ne correspondait pas — le cas que traite
cette section. Aucune de ces deux lignes n'enregistre jamais le mot de
passe tenté.
